\documentclass[12pt,a4paper,openany]{article}
\usepackage[english,russian]{babel}
\usepackage{bstyle}

\begin{document}
\section{Описание}
Программа SpcGen декодирует 16-значные значения из бинарных файлов с результатми оцифровки с АЦП и обрабатывает для получения спектра с детекторов.
В настоящий момент поддерживает два АЦП -- DRS4 и Digitizer. Также программа поддерживает два метода обработки сигнала -- по заряду и амплитуде.
Конфигурация программы реализована двумя парсерами -- с помощью конфигурационного файла и консольная.
Исходный код программы находится в репозитории: \href{https://github.com/kuzmenkoas/SpcGen}{GitHub}.

\section{Зависимости}
\begin{itemize}
    \item ROOT
    \item cmake and make
    \item C++17
\end{itemize}

\section{Инструкция по установке}
\subsection{Linux}
Установка зависимостей (для Ubuntu/Linux Mint):
\begin{lstlisting}[language=bash]
sudo apt install gcc cmake
\end{lstlisting}

Установка программы:
\begin{lstlisting}[language=bash]
git clone https://github.com/kuzmenkoas/SpcGen
cd SpcGen
mkdir build
cd build
cmake ..
make
\end{lstlisting}

\subsection{Windows}
\begin{enumerate}
    \item Установить \href{https://visualstudio.microsoft.com/downloads/}{Visual Studio}.
    \item Установить \href{https://cmake.org/download/}{CMake}.
    \item Собрать проект с помощью утилиты CMake (использовать компилятор MSVC).
    \item Скомпилировать проект с помощью Visual Studio (кнопка SpcGen.exe).
    \item Добавить путь к исполняемому файлу (путь/SpcGen/out/build/x64-Debug/) в переменную окружения PATH (Win+R, вызвать sysdm.cpl).
\end{enumerate}

\section{Запуск}
Для запуска необходимо указать путь к исполняемому файлу (в Linux build/SpcGen), в Windows, если добавить путь в переменную окружения PATH, то из любой директории можно вызвать SpcGen.exe.

Программа позволяет использовать два варианта задания конфигурации декодирования (с помощью командной строки и с помощью конфигурационного файла).

В директории /examples/data лежат бинарные файлы с DRS4 и Digitizer, а также в директории /examples/config лежат примеры конфигурационных файлов.

Программа в данный момент разделяет входные данные (аргументы) на конфигурационные (расширение .cfg) и бинарные (для drs -- .dat, для digitizer -- .bin), таким образом, 
что очередность передачи аргументов не влияет на программу (можно сначала передать конфигурационный, а далее бинарный, или сначала бинарные файлы, а после -- конфигурационный файл).

Определение типа устройства реализовано по расширению бинарного файла.

Для запуска с помощью конфигурационного файла устройства drs необходимо выполнить команду:
\begin{lstlisting}
    SpcGen.exe /examples/config/drs.cfg /examples/data/drs/08_Co60_ch4_NaI_N5_OC_R1307_1500V_30mV_T20_42_1Gss_500ns_20250618s.dat
\end{lstlisting}

Если использовать консольный парсер без конфигурационного файла, то:
\begin{lstlisting}
    SpcGen.exe /examples/data/drs/08_Co60_ch4_NaI_N5_OC_R1307_1500V_30mV_T20_42_1Gss_500ns_20250618s.dat
\end{lstlisting}

У устройства digitizer есть два типа бинарных файлов -- с данными обработки самой программой psd и с данными формы сигнала waveform.
Соответственно, есть три варианта декодирования бинарных файлов:
\begin{enumerate}
    \item psd;
    \item waveform;
    \item psd и waveform.
\end{enumerate}

В случае одновременной обработки psd и waveforn уже обязательна последовательность передачи бинарных файлов. 
Первым аргументом из бинарных файлов должен быть файл psd, следующим -- waveform.

Примеры обработки только psd файла:
\begin{lstlisting}
    SpcGen.exe /examples/config/psd.cfg /examples/data/digitizer/raw_data_psd_wavetest_2025_11_06__13_28_09.bin
\end{lstlisting}

\begin{lstlisting}
    SpcGen.exe /examples/data/digitizer/raw_data_psd_wavetest_2025_11_06__13_28_09.bin
\end{lstlisting}

Примеры обработки только waveform файла:
\begin{lstlisting}
    SpcGen.exe /examples/config/waveform.cfg /examples/data/digitizer/raw_waveform_psd_wavetest_2025_11_06__13_28_09.bin
\end{lstlisting}

\begin{lstlisting}
    SpcGen.exe /examples/data/digitizer/raw_waveform_psd_wavetest_2025_11_06__13_28_09.bin
\end{lstlisting}

Примеры обработки psd и waveform файлов:
\begin{lstlisting}
    SpcGen.exe /examples/config/psdWaveform.cfg /examples/data/digitizer/raw_data_psd_wavetest_2025_11_06__13_28_09.bin /examples/data/digitizer/raw_waveform_psd_wavetest_2025_11_06__13_28_09.bin
\end{lstlisting}

\begin{lstlisting}
    SpcGen.exe /examples/data/digitizer/raw_data_psd_wavetest_2025_11_06__13_28_09.bin /examples/data/digitizer/raw_waveform_psd_wavetest_2025_11_06__13_28_09.bin
\end{lstlisting}

\section{Конфигурационные файлы}
Общие заголовки для обоих АЦП -- Output, Signal, Histogram. Очередность заголовков не имеет значения. 
Программа реализована так, что ищет нужный заголовок и останавливается тогда, когда видит пустую строку.
Также, очередность параметров под заголовком не строго типизирована.

\subsection{Формат вывода}
Формат вывода настраивается под заголовком Output.
Далее через символ <<+>> через пробел пишется формат вывода. Поддерживается два вывода -- Root и Txt.
Пример:
\begin{lstlisting}
Output
+ Root
+ Txt
\end{lstlisting}

\subsection{Настройка обработки амплитудным методом}
Заголовком для амплитудного метода является Signal.
Далее параметрами являются: up (down) и range.
Параметр вида сигнала (Пик сигнала внизу (down), или наверху (up)) настраивает алгоритм так, чтобы искал минимальное (максимальное) значение как базовую точку для фитирования.
Далее, параметр range настраивает область фитирования, передается количество бинов слева и справа, где фитировать от базовой точки.
Пример:
\begin{lstlisting}
Signal
+ up
+ range 100 100
\end{lstlisting}

\subsection{Гистограммы}
Histogram является заголовком для настройки построения гистограмм по данным.
Далее в зависимости от устройства немного различается настройка.

Для АЦП DRS4 последовательность параметров:
\begin{enumerate}
    \item параметр;
    \item количество бинов;
    \item минимальное значение;
    \item максимальное значение.
\end{enumerate}

Для АЦП digitizer последовательность:
\begin{enumerate}
    \item файл (PSD или waveform);
    \item параметр;
    \item количество бинов;
    \item минмиальное значение;
    \item максимальное значение.
\end{enumerate}

Пример для АЦП DRS4:
\begin{lstlisting}
Histogram
+ charge 1000 -5 70
+ amplitude 1000 -0.05 0.6
\end{lstlisting}

Пример для АЦП Digitizer:
\begin{lstlisting}
Histogram
+ PSD qShort 1000 -20000 200000
+ Waveform charge 1000 -20000 200000
\end{lstlisting}

\subsection{Данные для обработки и декодирования}
Тут заголовочные файлы отличаются.

\end{document}